\section{The web target}
\label{sec:web}

\textsf{Y} compiles a second kind of executable beside the hosted ELF and
the bootable kernel: \ic{--target=web} writes a WebAssembly module for
\ic{wasm32}, and the same program runs under Node or in a page.  The target is
not an interpreter and not a bytecode: the compiler's own IR is lowered to
wasm, the context and the token encoding are fixed, and a JavaScript host
supplies the kernel calls the program makes.  ABI 1 names the module's
exports, its metadata section and the host calls; \ic{docs/TARGETS.md} is the
reference table.

\begin{lstlisting}
$ ./ycc --target=web prog.y -o prog.wasm
$ node --disable-wasm-trap-handler src/web/yhost.js run prog.wasm
$ ./ycc --target=web --web-page -o site prog.y   # index.html + prog.wasm
$ ./ycc build --target=web                        # y.toml: build/web/NAME/
\end{lstlisting}

\paragraph{The module.} Values are 64-bit; memory is wasm32 and grows on
demand.  Each Y function body is a \emph{chunk}: a loop with a \ic{br\_table}
on a token and \ic{return\_call}/\ic{return\_call\_indirect} between chunks,
so a call does not grow the wasm operand stack and 1000-deep recursion costs
no host stack.  A token is \ic{2\string^40 + (chunk << 20) + case}, a
64-bit word whose low bits are the case and whose high bits say which
function owns the continuation.  The module exports \ic{memory},
\ic{resume(i64)}, \ic{ctx\_save(i32)}, \ic{ctx\_load(i32)},
\ic{set\_fuel(i32)}, \ic{why()} and \ic{set\_why(i32)}; \ic{resume}
runs until it returns a token, and the host turns that into the next
\ic{resume}, a syscall or a trap.  The custom section \ic{yweb} carries
\ic{abi}, \ic{stack\_lo}, \ic{stack\_hi}, \ic{data\_lo}, \ic{data\_end},
\ic{heap\_lo}, \ic{entry}, \ic{chunks}, \ic{threads}, \ic{gui} and
\ic{audio}.  The initial stack is the Linux one (argc, argv, envp, auxv), so
\ic{argv()} and the environment work unchanged.

\paragraph{The host.} \ic{src/web/core.js} is platform-neutral: it loads and
validates a module, owns the linear memory, the green-thread scheduler, the
syscall table and the host-call table.  \ic{vfs.js} is one VFS for both hosts
(descriptors, pipes, eventfd/timerfd, \ic{/proc}, \ic{/dev}); the Node backend
is the real filesystem, the browser backend an in-memory tree with
\ic{/app} from \ic{files.bin} and \ic{/home/web} in IndexedDB.  Threads are
green threads with a fuel budget: a slice ends after its fuel, a blocking
call parks the thread and the scheduler resumes the next one, and
\ic{pause\_hidden} freezes a hidden page's clock and scheduler.  A raw Linux
syscall is answered by the VFS when it has a plausible answer and refused
with the errno the program would get from a kernel that does not have the
feature; \ic{fork}, \ic{exec}, sockets and \ic{obfuscate} are refused by
design, each with its reason in \ic{test/y/native/*.skip.web}.

\paragraph{The host API.} \ic{y\_os\_host\_sym(name, len)} returns a
positive index for a host function and 0 when there is none;
\ic{y\_os\_host\_call(h, args, n)} calls it with \ic{n} i64 arguments, and a
Promise blocks the calling green thread until it settles.  Strings come back
through a \ic{(buf, cap)} pair: the function returns the length, or
\ic{-(needed)} when the buffer is too small.  The built-in names are
\ic{gui.*}, \ic{audio.*} and \ic{web.*}; a page may register its own under
\ic{app.*}.  \ic{std/web} wraps the page's own API: \ic{available},
\ic{platform}, \ic{sym}, \ic{call}, \ic{call\_str}, \ic{eval}, \ic{log},
\ic{param}, \ic{open\_url} and \ic{fetch}; off the web every one answers
\ic{false}, \ic{""}, \ic{0} or \ic{none}, so the same source still builds and
runs on Linux and Windows.

\paragraph{The screen and sound.} \ic{std/gui} keeps one API.  For web it is
\ic{gui.web.y}, a target overlay merged over \ic{gui.y} and packed as
\ic{gui@web}: the nine \ic{gui.*} calls go through the host symbols, while
present/pump/wait/resize/read\_back/verify\_present/close keep the X11 and
Win32 arms for native builds.  The page owns a canvas: fixed or fit sizing,
integer or automatic zoom (default 1, 100\%), a pixel buffer swapped from
B G R A, and the full
key, char, pointer, wheel, focus, composition and paste input map, with
\ic{ResizeObserver} and \ic{pagehide} becoming events.  Sound is the four
\ic{y\_os\_audio\_*} calls: the browser plays through an inline
\ic{AudioWorklet} (a \ic{ScriptProcessorNode} fallback when a CSP refuses the
Blob URL) and parks a thread while more than 250 ms is queued; Node writes a
WAV with \ic{--audio=FILE.wav}.  \ic{web.available()} is the feature test a
program uses before it asks for a window or a speaker.

\paragraph{Page, site and deployment.} The page provides the stage the
program runs in: one full-window area.  A program that opens a gui window
draws into the viewport, and one without a window gets the console (the page's
stdin/stdout/stderr); \ic{[web] console} = \ic{"show"} or \ic{"hide"} chooses.
A \ic{[web]} table in \ic{y.toml} sets the
title, id, start mode, canvas mode, zoom (100\% unless the app asks otherwise),
console, persistence, \ic{files.bin}, arguments, environment, stack,
pause\_hidden and the single-file form; a \ic{template} points the generator
at the application's own page, whose four placeholders are \ic{TITLE},
\ic{CONFIG}, \ic{HOST} and \ic{MODULE}.  See \ic{docs/TARGETS.md}.

\ic{--web-page DIR} writes \ic{index.html} (the stylesheet and the linked host
are inline) and \ic{NAME.wasm}, plus \ic{files.bin} for \ic{[web] files};
the deployment is those files.  \ic{--web-page=html} additionally inlines the
module and the bundle, so it is one file.  \ic{ycc build --target=web} writes
the same site under \ic{build/web/NAME/}, a single file under
\ic{build/web/NAME.html}, and a Node launcher under \ic{build/web/bin/NAME}
with its private host copy.  A module is fetched as a sibling or as a data URL.
\ic{?ytest=1} writes \ic{YDONE} into the title and \ic{<pre id="ytest">},
which is how the browser suite drives a real page; \ic{?arg=} passes a program
argument and \ic{?fuel=} a slice budget.  The page is inline JavaScript, so a
deployment with a strict CSP must allow it (the module is still fetched, not
eval'd).

\paragraph{What runs.} The whole \ic{test/y/native} corpus runs through the
Node host, and the fixtures that need a system the web does not have skip
with a reason: \ic{fork}, sockets, \ic{exec}, and obfuscation.  Threads,
sleep order, preemption, the collector, NaN and division flags, the token
burst, big bodies, traps, checked builds, stdin, the page API, gui and sound
all have web fixtures.
Performance and size are measured, not guessed: the wasm module is several
times a native build's work but a browser runs it inside a frame, and the
network cost is the module's gzip size.

\paragraph{Limits.} A web build has no \ic{fork}, \ic{exec}, \ic{popen},
sockets, \ic{/proc/self/mem} or \ic{ptrace}; a program that needs a real
kernel stays native.  \ic{--pack}, \ic{--obfuscate}, \ic{--target=kernel} and
\ic{-g} either refuse the target or are not part of it.  The host is not a
security boundary: a page's JavaScript can read its own module and memory as
any page can, and a save in IndexedDB is the player's.
